\section{Conclusion}\label{sec:conclusion}

We have established the first rigorous algebraic framework for chain-of-thought
prompt composition, modeling CoT operators as elements of a semigroup acting on
a metric space of reasoning trajectories. The five main results of this
paper---the monoid structure (Theorem~\ref{thm:semigroup}), the sharp
contractivity criterion (Theorem~\ref{thm:contractivity}), the
idempotent--fixed point correspondence (Theorem~\ref{thm:idempotent}), the
four-case convergence classification (Theorem~\ref{thm:classification}), and
the JdLG decomposition (Theorem~\ref{thm:jdlg})---provide a complete algebraic
characterization of when and how iterative CoT prompting converges.

The central insight is that convergence of iterated CoT application is governed
by a single computable quantity: the spectral radius $\rho(B)$ of the
current-state transformation matrix. When $\rho(B)<1$, iterates converge
exponentially to a consensus trajectory determined by the initial query; when
$\rho(B)=1$, the system oscillates or stabilizes depending on the spectral
structure of~$B$; and when $\rho(B)>1$, the system diverges. The JdLG
decomposition further separates the trajectory space into stable reasoning
patterns and transient noise, providing a principled notion of
``signal versus noise'' in reasoning trajectories.

This framework opens several directions for future investigation, including
extension to nonlinear and stochastic operators, computation of joint spectral
radii for heterogeneous prompt sequences, and empirical validation in
high-dimensional embedding spaces. We anticipate that the algebraic perspective
developed here will provide useful tools for the principled design and analysis
of iterative reasoning systems.
